\begin{theorem}[6.12.7]
\label{IV.2.s6.12.7-fr}
\textnormal{(Nagata).} --- Sean $A$ un anillo local noetheriano completo y
$X=\operatorname{Spec}(A)$. Entonces $\operatorname{Reg}(X)$ es abierto en
$X$.
\end{theorem}

\begin{proof}
Teniendo en cuenta (6.12.2), la cuestión se reduce al caso en que, además,
$A$ es \emph{íntegro}, y a demostrar que en este caso $\operatorname{Reg}(X)$
contiene una parte abierta no vacía de $X$. Distingamos dos casos:

\oldpage[IV-2]{167}
I) \emph{El cuerpo de fracciones $K$ de $A$ es de característica} 0. --- Se sabe
entonces (\textbf 0, 19.8.8) que existen un anillo de valuación discreta completo
$C$ y un subanillo $B$ de $A$ tales que $A$ es una $B$-álgebra \emph{finita} y
$B$ es isomorfo a un anillo de series formales
$C[[T_1,\ldots,T_r]]$. Como $C$ es regular (\textbf{II}, 7.1.6), también lo es
$B$ (\textbf 0, 17.3.8); además, como el cuerpo de fracciones $L$ de $B$ es de
característica 0, $K$ es una extensión \emph{separable} finita de $L$, y se puede
aplicar (6.12.4.1) a $B$, lo que demuestra la proposición.

II) \emph{El cuerpo de fracciones $K$ de $A$ es de característica} $p>0$. ---
Entonces $A$ contiene el cuerpo primo $\mathbf F_p$, y el teorema ya ha
sido demostrado en este caso (\textbf 0, 22.7.6).
\end{proof}
